\section{Conclusiones y trabajo futuro}
A partir de la metodología cualitativa implementada, se evidenció que la combinación de una arquitectura sólida (MVC/REST) y una metodología ágil (Scrum) resulta fundamental para el éxito en el desarrollo de sistemas.
La revisión documental comparativa entre Spring Boot y Node.js confirmó que no existe un framework superior, sino que la elección se basa en la correcta alineación entre la tecnología seleccionada, los objetivos del proyecto y las competencias del equipo:
•	Spring Boot: Ideal para proyectos que requieren solidez, alta seguridad y madurez (entornos empresariales).
•	Node.js: Mejor opción para aplicaciones que necesitan flexibilidad, rapidez y alta escalabilidad con tiempos de respuesta cortos.
La aplicación del patrón MVC facilita la organización, escalabilidad y mantenibilidad del código, mientras que el uso de Scrum fortaleció la planificación y la entrega incremental de funcionalidades. Estos hallazgos validan que la investigación teórica y la práctica de desarrollo backend deben estar alineadas para lograr un producto funcional y sustentable a largo plazo.